\section{Related work, limitations and outlook}\label{sec:related}

\subsection{QAM and proof assistants}

A proof assistant such as Lean~4~\cite{lean4} checks the \emph{content} of mathematical statements: a theorem formalised in Lean with
Mathlib~\cite{mathlib} is proved, relative to the kernel of the system. QAM checks the \emph{form} of a record: which status a statement
may have given the status of its premises, which inputs it carries and in which frame it holds. The two answer different questions and do
not exclude each other.
\begin{center}
\small
\begin{tabular}{@{}p{0.25\textwidth}p{0.33\textwidth}p{0.33\textwidth}@{}}
\toprule
 & Proof assistant & QAM\\
\midrule
checks & validity of a proof & consistency of status, inputs, frames\\
cost of entry & full formalisation & one declaration per statement\\
open statements & \texttt{sorry} or an axiom & a first-class status, inherited\\
physical interpretations & outside the system & recordable as \st{reading}\\
calculations & inside the system or as an oracle & as a request with a receipt\\
\bottomrule
\end{tabular}
\end{center}
QAM uses Lean for its own checker and for the abstract machine. It does not formalise the content of records. A natural extension is to
link a Lean proof to an entry with status \st{proved}, so that this status is evidenced rather than declared. No such link exists yet.

\subsection{Related approaches}

\paragraph{Dependency graphs of mathematical texts.} Closest are the \emph{blueprints} of large Lean formalisation projects. The tool
\texttt{leanblueprint}~\cite{blueprint} reads a \LaTeX{} text in which theorems and definitions state their dependencies, builds a
dependency graph and tracks the formalisation state of each node. KnowTeX~\cite{knowtex} extracts dependency graphs from annotated
\LaTeX{} sources, and LeanMarathon~\cite{leanmarathon} uses an evolving blueprint as scaffolding for formalisation by several
language-model agents; its diagnosis (statements shift, dependencies tangle, context is lost) agrees with ours. The difference lies in the
status. There it measures progress towards formalisation, and the goal is a complete proof in the assistant. In QAM it measures the
epistemic strength of a statement, including interpretations that cannot be formalised, and a status above that of the premises is an
error.

\paragraph{Language models and formal mathematics.} Autoformalisation translates informal mathematics into the language of a proof
assistant~\cite{autoformal}, and language models are used as proof search engines in Lean~\cite{leandojo}. Both answer
hallucination~\cite{halluc} with a complete checker for content. QAM starts one step earlier, with statements that are not (yet)
formalisable, and checks only their structure.

\paragraph{Assumption management.} Truth maintenance systems keep a justification for every inference~\cite{doyle79}; de Kleer's
assumption-based system labels every inference with the sets of assumptions under which it holds~\cite{atms}. QAM's input sets are a
restricted form of such labels, without the management of contradictions. In safety cases written in the Goal Structuring
Notation~\cite{gsn}, claims are linked to evidence, assumptions and explicitly undeveloped goals; an undeveloped goal plays the role of
a guard.

\paragraph{Lattices and provenance.} The status rule is a lattice model in the sense of Denning's information flow~\cite{denning76},
read in the opposite direction: instead of confidentiality flowing upwards, certainty can only flow downwards. The propagation of inputs
corresponds to provenance tracking in databases, as captured in general by provenance semirings~\cite{semiring}. The executor's log is a
provenance record in the sense of the W3C PROV model~\cite{prov}. Tools such as noWorkflow~\cite{noworkflow} capture the provenance of
script executions, and workflow languages such as CWL~\cite{cwl} and Snakemake~\cite{snakemake} describe computations reproducibly. They
record \emph{how} a result came about; QAM adds \emph{as what} it may be cited.

\paragraph{Types and proof-carrying code.} The properties proved for the machine have the form of syntactic type
soundness~\cite{wright94,pierce}; the rule that a certificate contains its proof is the idea of proof-carrying code~\cite{necula}.

\subsection{Limitations}

QAM checks only what is declared. A statement declared \st{proved} without premises is accepted; a fitted value declared again under a
new name is not recognised as an input; a false premise stays false. Frames are names, maps carry no status, and formulas are text. The
parser of the Lean checker is not verified, and neither is the running executor; MD5 protects against accidental change only. The tape
comparison covers a finite set of cases of one instruction. And the system has so far been used in one project, with one author: there
is no controlled study of whether it reduces errors, only the experience that it makes certain errors visible.

\subsection{Outlook: a quantum computer in the physics itself?}

\emph{This subsection is speculation. Nothing in it is established, and QAM would record every claim in it as \st{open}.}

If a mathematical structure can carry a classical tape (Section~\ref{sec:tape}), one may ask whether it could carry a quantum computer:
a machine whose tape holds superpositions of configurations and whose steps act on them coherently. The question is natural for a
carrier built from spinors, because a two-component spinor is exactly the state space of a qubit. It also has a long history: Feynman
proposed simulating physics with quantum systems~\cite{feynman82}, and Deutsch defined a universal quantum computer as a physical
system obeying quantum mechanics~\cite{deutsch85}.

What exists today is modest. The tape machine is classical and deterministic. Its Qiskit version builds each step as a \emph{reversible}
circuit, and reversibility is a precondition for coherent execution, since unitary evolution is reversible and an irreversible step
would have to dissipate information~\cite{bennett73}. A check confirms that the circuit followed by its inverse restores a superposition of two inputs.
But in the machine runs the circuit is applied to basis states only: every step is read out and the result
fed into the next one, so no superposition survives from step to step. The spinor tape hypothesis, as stated, concerns classical reading
and writing along the carrier.

At least four further conditions would be needed, each of which would be an open obligation in a QAM record:
\begin{enumerate}[label=\arabic*.,nosep,leftmargin=2em]
\item the carrier admits superpositions of distinct tape configurations as states, not only definite ones;
\item each machine step acts on those states as a unitary map, including head movement and control;
\item no step implicitly reads out the state: branch predicates must be applied coherently, as controlled operations, not as
measurements;
\item the coherence survives long enough, or can be protected well enough, for a computation of useful length.
\end{enumerate}
The third condition conflicts directly with the present design, in which branch predicates are \emph{read} (Section~\ref{sec:tape}). A
quantum version would need a different branching primitive. The fourth is the hardest problem of experimental quantum computing and has
no reason to be easier in a mathematical carrier that is meant to describe physics.

What QAM can contribute here is bookkeeping, not an answer. It would hold ``the carrier supports a quantum computer'' as an \st{open}
entry with these four conditions, separate from the conditional classical claim and from the executable checks of the classical plan, and
it would prevent any successful classical run from being cited as evidence for the quantum claim.

\subsection{Conclusion}

QAM does one thing reliably: it makes visible and checkable what each statement rests on, and it prevents an open condition from being
lost along a long chain of reasoning. Zuse's separation of plan and execution supplies the architecture: the language model calculates and
formulates, and the record decides what holds. Turing's boundary supplies the reason for the third outcome: what is not decided is not
smoothed away but kept as a guard, stored and inherited. Securing the form of a record does not make its content true. But it ensures
that one can see which part of the content is not yet secured.
